\subsection*{The theory of sets}

It may be said that at every period of history mathematicians and philosophers have used arguments of the theory of sets more or less consciously. But in the history of their conceptions on this subject, it is necessary to separate clearly all questions relating to the idea of cardinal numbers (and, in particular, relating to the notion of infinity) from those which involve only the notions of membership and inclusion. The latter are intuitively clear and seem never to have given rise to controversies. They serve as the easiest foundation for a theory of the syllogisms (as Leibniz and Euler were to show), for axioms such as ``the whole is greater than the part'', and for that part of geometry which is concerned with intersections of curves and surfaces. Up to the end of the nineteenth century, there was equally no objection raised to speaking of the set (or ``class'' according to some writers) of all objects having one or another given property (*); and the celebrated ``definition'' given by Cantor (``By a set we mean a grouping into one entity of distinct objects of our intuition or our thought'') ({[}25{]}, p.~282) provoked hardly any objections at the time of its publication (†). But the situation was quite different where the notions of number or magnitude became entangled with the notion of set. The question of the indefinite divisibility of length (raised, certainly, by the early Pythagoreans) led, as is well known, to considerable philosophical difficulties; from the Eleatic school to Bolzano and Cantor, mathematicians and philosophers unsuccessfully came up against the paradox of the finite magnitude composed of an infinite number of points without magnitude. For us it is irrelevant to retrace, even in summary, the interminable and passionate polemics aroused by this problem, which constituted a particularly favourable terrain for metaphysical or theological aberrations. Let us only take note of the point of view held since antiquity by the majority of mathematicians. It consisted essentially in refusing to discuss the question since it could not be irrefutably decided, an attitude which we shall find again among the modern formalists: just as the latter contrive to eliminate all appearances of ``paradoxical'' sets (see later, page 328), so the classical mathematicians carefully avoided introducing into their arguments the ``actual infinite'' (that is to say, sets containing an infinity of objects conceived as existing simultaneously, at least in thought) and contented themselves with the ``potential infinite'', i.e., the possibility of increasing any given magnitude (or of diminishing it, if a ``continuous'' magnitude was under consideration) (**). If this point of view contained a certain measure of hypocrisy (††), it nevertheless allowed the development of the greater part of classical mathematics (including the theory of proportion and, later, the infinitesimal calculus) (*); it appeared to be an excellent lifeline, especially after the quarrels aroused by infinitesimals, and had become an almost universally accepted dogma until well into the nineteenth century.

A first glimmer of the general notion of equipotence appeared in a remark of Galileo ({[}8{]}, vol.~VIII, pp.~78-80). He observed that the mapping \(n \to n^2\) established a one-to-one correspondence between the natural integers and their squares, and consequently that the axiom ``the whole is greater than the part'' could not be applied to infinite sets. But, far from inaugurating a rational study of infinite sets, this remark seems to have had no other effect than to reinforce distrust of the ``actual infinite''; this was the moral that Galileo himself drew, and Cauchy quotes him in 1833 only to endorse his attitude.

The requirements of analysis --- and especially the deeper study of functions of real variables, which was pursued throughout the nineteenth century --- are at the origin of what was to become the modern theory of sets. When Bolzano in 1817 proved the existence of the greatest lower bound of a set bounded below in R, he argued, as most of his contemporaries did, in terms of ``comprehension'' and considered, not an arbitrary set of real numbers, but an arbitrary property of them. But when, thirty years later, he wrote his Paradoxien des Unendlichen {[}18{]} (published in 1851, three years after his death), he did not hesitate to claim the possibility of existence for the ``actual infinite'' and to speak of arbitrary sets. In this work he defined the general notion of equipotence of two sets and proved that any two compact intervals in R are equipotent (provided each contains more than a single point). He observed also that the characteristic difference between finite and infinite sets lies in the fact that an infinite set E is equipotent to some proper subset of E, but he gave no convincing proof of this assertion. The general tone of the work is more philosophical than mathematical; and since he did not distinguish precisely enough between the notion of the power of a set and that of magnitude or order of infinity, Bolzano failed in his attempts to construct infinite sets of greater and greater powers, and his discussion in this connection was entangled with considerations of divergent series which are entirely devoid of meaning.

The creation of the theory of sets, as it is understood today, is due to the genius of G. Cantor. He too began as an analyst, and his work on trigonometric series, inspired by that of Riemann, led him naturally in 1872 to a first essay in classifying the ``exceptional'' sets which arise in this theory (*) through the notion of successive ``derived sets'', which he introduced in this connection. Undoubtedly it was these researches, and also his method for defining real numbers, which led Cantor to interest himself in problems of equipotence; for in 1873 he noted that the set of rational numbers (or the set of algebraic numbers) is countable. In his correspondence with Dedekind, which began at about this time {[}25 bis{]}, he raised the question whether the set of integers and the set of real numbers are equipotent, and succeeded a few weeks later in showing that they are not. Next, from 1874 onward, the problem of dimension preoccupied him, and for three years he sought in vain to establish the impossibility of a one-to-one correspondence between R and \(R^{n}\) (n \textgreater{} 1), before he came to construct such a correspondence, to his own stupefaction (†). Having obtained these new and surprising results, he devoted himself entirely to the theory of sets. In a series of six memoirs published in the Mathematische Annalen between 1878 and 1884, he considered problems of equipotence, the theory of totally ordered sets, topological properties of R and \(R^{n}\) , and the problem of measure; and it is admirable to see how in his hands the notions which appeared to be so inextricably rooted in the classical conception of the ``continuum'' were gradually brought out. In 1880 he had the idea of iterating ``transfinitely'' the formation of ``derived sets''; but this idea did not take substance until two years later, with the introduction of well-ordered sets --- one of the most original of Cantor's discoveries, thanks to which he was able to initiate a detailed study of cardinal numbers and to formulate the ``problem of the continuum'' {[}25{]}.

It was not to be expected that such bold conceptions, which ran counter to a tradition two thousand years old and led to such unlikely results of so paradoxical an appearance, should have been accepted without resistance. In fact, among the influential mathematicians in Germany, Weierstrass was the only one to follow with some favour the work of Cantor (his former pupil); but Cantor met with the irreconcilable opposition of Schwarz, and above all of Kronecker (*). It seems to have been as much the constant tension generated by opposition to his ideas as his fruitless efforts to prove the continuum hypothesis which produced in Cantor the first symptoms of a nervous illness which was to affect his mathematical productivity (†). In fact he did not resume his interest in the theory of sets until about 1887, and his last publications, which are concerned mainly with the theory of totally ordered sets and the calculus of ordinals, date from the period 1895-97. He had also proved in 1890 the inequality m \textless{} 2 \(^{m}\) . However, not only did the problem of the continuum remain unanswered, but there was a more serious gap in the theory of cardinals, for Cantor had not been able to prove the existence of a well-ordering relation between arbitrary cardinals. This gap was to be filled partly by the theorem of F. Bernstein (1897) showing that the relations a ≤ b and b ≤ a imply a = b (**), and above all by the theorem of Zermelo {[}36 a{]} which proved the existence of a well-ordering on any set: a theorem that Cantor had conjectured already in 1883 ({[}25{]}, p.~169).

Dedekind, however, had followed the work of Cantor with sustained interest since the beginning. But whereas the latter concentrated his attention on infinite sets and their classification, Dedekind pursued his own thoughts on the notion of number (which had already led to his definition of irrational numbers by means of ``cuts''). In his pamphlet Was sind und was sollen die Zahlen, which was published in 1888, but whose essential content dates from the period 1872-8 ({[}26{]}, vol.~III, p.~335), he showed how the notion of natural integer (which, as we have seen, was now the basis for the whole of classical mathematics) could itself be derived from the fundamental notions of the theory of sets. He was certainly the first to develop explicitly the elementary properties of arbitrary mappings of one set into another (hitherto neglected by Cantor, who was interested only in injective correspondences) and introduced, for any mapping f of a set E into itself, the notion of the ``chain'' of an element \(a \in E\) relative to f, namely the intersection of all sets \(K \subset E\) such that \(a \in K\) and \(f(K) \subset K (*)\) . Next, he took as the definition of an infinite set E the fact that there exists a one-to-one mapping \(\varphi\) of E into E such that \(\varphi(E) \neq E (\dagger)\) . If, furthermore, there exists such a mapping \(\varphi\) and an element \(a \notin \varphi(E)\) for which the whole set E is the chain of a, Dedekind said that E is ``simply infinite''; moreover, he noted that ``Peano's axioms'' are then satisfied and showed (before Peano did) how, from that starting-point, all the elementary theorems of arithmetic may be obtained. The only thing lacking in his presentation is the axiom of infinity, which Dedekind (following Bolzano) believed he could prove by considering the ``world of thoughts'' (Gedankenwelt) as a set (**).

On another front, Dedekind had been led by his arithmetical work (and especially the theory of ideals) to envisage the notion of ordered set in a more general aspect than Cantor had done. Whereas the latter restricted himself entirely to totally ordered sets (*), Dedekind investigated the general case and in particular made a close study of lattices ({[}26{]}, vol.~II, pp.~236-271). This work attracted hardly any attention at the time; and although his results, rediscovered by various authors, have been the object of numerous publications since 1935, their historical importance lies far less in the possible applications of this theory (which indeed are rather scanty) than in the fact it was one of the first examples of careful axiomatic construction. On the other hand, Cantor's first results on countable sets soon produced many important applications, even to the most classical questions of analysis (†) (quite apart from those parts of his work which inaugurated general topology and measure theory; for these the reader is referred to the Historical Notes in General Topology). Moreover, the last years of the nineteenth century saw the first applications of the principle of transfinite induction which, especially since the proof of Zermelo's theorem, has become an indispensable instrument in all parts of modern mathematics. Kuratowski in 1922 gave a version of this principle, subsequently rediscovered by Zorn {[}45{]} and customarily named after him, which is often more convenient to apply since it avoids the use of well-ordered sets ({[}40{]}, p.~89); it is this form of the principle that is generally used nowadays (**).

By the end of the nineteenth century, therefore, the essential conceptions of Cantor had carried the day (††). As we have seen, the formalization of mathematics was achieved at the same period, and the use of the axiomatic method had become almost universally accepted. In other words, the intense labours of the years 1875-1895 had acquired for mathematics the essential content of the material presented in this book. But the same period saw the beginning of a ``crisis of the foundations'' of rare violence, which was to shake the mathematical world for more than thirty years and seemed at times to put in danger not only all the recent acquisitions but even the most classical parts of mathematics.

